
\documentclass{article}
\usepackage{colortbl}
\usepackage{makecell}
\usepackage{multirow}
\usepackage{supertabular}

\begin{document}

\newcounter{utterance}

\centering \large Interaction Transcript for game `dond', experiment `coop', episode 15 with Qwen3.6{-}35B{-}A3B{-}FP8{-}without{-}reasoning.
\vspace{24pt}

{ \footnotesize  \setcounter{utterance}{1}
\setlength{\tabcolsep}{0pt}
\begin{supertabular}{c@{$\;$}|p{.15\linewidth}@{}p{.15\linewidth}p{.15\linewidth}p{.15\linewidth}p{.15\linewidth}p{.15\linewidth}}
    \# & $\;$A & \multicolumn{4}{c}{Game Master} & $\;\:$B\\
    \hline

    \theutterance \stepcounter{utterance}  
    & & \multicolumn{4}{p{0.6\linewidth}}{
        \cellcolor[rgb]{0.9,0.9,0.9}{
            \makecell[{{p{\linewidth}}}]{
                \texttt{\tiny{[P1$\langle$GM]}}
                \texttt{Sie spielen ein kooperatives Verhandlungsspiel, bei dem Sie sich mit einem anderen Spieler darauf einigen müssen, wie eine Reihe von Gegenständen aufgeteilt werden soll.} \\
\\ 
\texttt{Die Regeln:} \\
\texttt{(a) Sie und der andere Spieler erhalten eine Sammlung von Gegenständen. Jeder von Ihnen erhält außerdem eine geheime Wertfunktion, die angibt, wie viel Ihnen jede Art von Gegenstand wert ist.} \\
\texttt{(b) Sie tauschen Nachrichten mit dem anderen Spieler aus, um zu vereinbaren, wer welche Gegenstände bekommt. Sie können jeweils maximal 5 Nachrichten senden oder das Spiel vorzeitig beenden, indem Sie jederzeit einen geheimen Vorschlag machen.} \\
\texttt{(c) Jeder von euch wird aufgefordert, einen geheimen Vorschlag zu machen, in dem ihr die gewünschten Gegenstände in eckigen Klammern wie folgt angibt: "{[}Vorschlag: <Nummer> <Objektname>, <Nummer> <Objektname>, <...>{]}"} \\
\texttt{(d) Wenn eure Vorschläge komplementär sind, d.h. es gibt genug Gegenstände, um beide Vorschläge zu erfüllen, erhält jeder Spieler eine Punktzahl, die sich aus der Summe der Werte für die Gegenstände ergibt, die er erhalten hat. Andernfalls erhalten Sie beide null Punkte.} \\
\texttt{(e) Ihr Ziel ist es, die Summe aus Ihrer Punktzahl und der Punktzahl des anderen Spielers zu maximieren.} \\
\\ 
\texttt{Beginnen wir.} \\
\\ 
\texttt{Die Menge der verfügbaren Gegenstände ist:} \\
\\ 
\texttt{3 Statuen, 1 Nagel, 1 Hammer} \\
\\ 
\texttt{Deine geheime Wertfunktion ist:} \\
\\ 
\texttt{Statue: 3, Nagel: 1, Hammer: 0} \\
\\ 
\texttt{WICHTIG: Ihre Nachrichten werden, sofern es sich nicht um einen geheimen Vorschlag handelt, direkt an den anderen Spieler übermittelt, also fügen Sie keine Antwort auf die Regeln oder einen Text zur Ankündigung Ihrer Nachricht ein. Um einen geheimen Vorschlag zu machen, verwenden Sie das angegebene Format. Verwenden Sie keine eckigen Klammern, wenn Sie mit dem anderen Spieler kommunizieren, sonst wird dies als Ihr geheimer Vorschlag interpretiert.} \\
            }
        }
    }
    & \\ \\

    \theutterance \stepcounter{utterance}  
    & \multicolumn{4}{p{0.6\linewidth}}{
        \cellcolor[rgb]{0.9,0.9,0.9}{
            \makecell[{{p{\linewidth}}}]{
                \texttt{\tiny{[P1$\rangle$GM]}}
                \texttt{Hallo! Ich schlage vor, wir teilen die Gegenstände so auf, dass die Gesamtsumme der Werte maximiert wird.} \\
\\ 
\texttt{Da ich den Hammer nicht brauche (Wert 0) und die Statue für mich wertvoller ist (3) als der Nagel (1), schlage ich folgende Aufteilung vor:} \\
\\ 
\texttt{Ich nehme: 2 Statuen, 1 Nagel} \\
\texttt{Du bekommst: 1 Statue, 1 Hammer} \\
\\ 
\texttt{So erhalte ich einen Wert von 3*2 + 1 = 7.} \\
\texttt{Du erhältst einen Wert von 3*1 + 0 = 3.} \\
\texttt{Gesamtwert: 10.} \\
\\ 
\texttt{Wenn du einverstanden bist, antworte bitte mit "Einverstanden" oder mache deinen eigenen Vorschlag.} \\
            }
        }
    }
    & & \\ \\

    \theutterance \stepcounter{utterance}  
    & & & \multicolumn{2}{p{0.3\linewidth}}{
        \cellcolor[rgb]{0.9,0.9,0.9}{
            \makecell[{{p{\linewidth}}}]{
                \texttt{\tiny{[GM$|$GM]}}
                \texttt{continue} \\
            }
        }
    }
    & & \\ \\

    \theutterance \stepcounter{utterance}  
    & & & \multicolumn{2}{p{0.3\linewidth}}{
        \cellcolor[rgb]{0.9,0.9,0.9}{
            \makecell[{{p{\linewidth}}}]{
                \texttt{\tiny{[GM$|$GM]}}
                \texttt{Hallo! Ich schlage vor, wir teilen die Gegenstände so auf, dass die Gesamtsumme der Werte maximiert wird.} \\
\\ 
\texttt{Da ich den Hammer nicht brauche (Wert 0) und die Statue für mich wertvoller ist (3) als der Nagel (1), schlage ich folgende Aufteilung vor:} \\
\\ 
\texttt{Ich nehme: 2 Statuen, 1 Nagel} \\
\texttt{Du bekommst: 1 Statue, 1 Hammer} \\
\\ 
\texttt{So erhalte ich einen Wert von 3*2 + 1 = 7.} \\
\texttt{Du erhältst einen Wert von 3*1 + 0 = 3.} \\
\texttt{Gesamtwert: 10.} \\
\\ 
\texttt{Wenn du einverstanden bist, antworte bitte mit "Einverstanden" oder mache deinen eigenen Vorschlag.} \\
            }
        }
    }
    & & \\ \\

    \theutterance \stepcounter{utterance}  
    & & \multicolumn{4}{p{0.6\linewidth}}{
        \cellcolor[rgb]{0.9,0.9,0.9}{
            \makecell[{{p{\linewidth}}}]{
                \texttt{\tiny{[GM$\rangle$P2]}}
                \texttt{Sie spielen ein kooperatives Verhandlungsspiel, bei dem Sie sich mit einem anderen Spieler darauf einigen müssen, wie eine Reihe von Gegenständen aufgeteilt werden soll.} \\
\\ 
\texttt{Die Regeln:} \\
\texttt{(a) Sie und der andere Spieler erhalten eine Sammlung von Gegenständen. Jeder von Ihnen erhält außerdem eine geheime Wertfunktion, die angibt, wie viel Ihnen jede Art von Gegenstand wert ist.} \\
\texttt{(b) Sie tauschen Nachrichten mit dem anderen Spieler aus, um zu vereinbaren, wer welche Gegenstände bekommt. Sie können jeweils maximal 5 Nachrichten senden oder das Spiel vorzeitig beenden, indem Sie jederzeit einen geheimen Vorschlag machen.} \\
\texttt{(c) Jeder von euch wird aufgefordert, einen geheimen Vorschlag zu machen, in dem ihr die gewünschten Gegenstände in eckigen Klammern wie folgt angibt: "{[}Vorschlag: <Nummer> <Objektname>, <Nummer> <Objektname>, <...>{]}"} \\
\texttt{(d) Wenn eure Vorschläge komplementär sind, d.h. es gibt genug Gegenstände, um beide Vorschläge zu erfüllen, erhält jeder Spieler eine Punktzahl, die sich aus der Summe der Werte für die Gegenstände ergibt, die er erhalten hat. Andernfalls erhalten Sie beide null Punkte.} \\
\texttt{(e) Ihr Ziel ist es, die Summe aus Ihrer Punktzahl und der Punktzahl des anderen Spielers zu maximieren.} \\
\\ 
\texttt{Beginnen wir.} \\
\\ 
\texttt{Die Menge der verfügbaren Gegenstände ist:} \\
\\ 
\texttt{3 Statuen, 1 Nagel, 1 Hammer} \\
\\ 
\texttt{Deine geheime Wertfunktion ist:} \\
\\ 
\texttt{Statue: 0, Nagel: 1, Hammer: 9} \\
\\ 
\texttt{WICHTIG: Ihre Nachrichten werden, sofern es sich nicht um einen geheimen Vorschlag handelt, direkt an den anderen Spieler übermittelt, also fügen Sie keine Antwort auf die Regeln oder einen Text zur Ankündigung Ihrer Nachricht ein. Um einen geheimen Vorschlag zu machen, verwenden Sie das angegebene Format. Verwenden Sie keine eckigen Klammern, wenn Sie mit dem anderen Spieler kommunizieren, sonst wird dies als Ihr geheimer Vorschlag interpretiert.} \\
\\ 
\texttt{Hallo! Ich schlage vor, wir teilen die Gegenstände so auf, dass die Gesamtsumme der Werte maximiert wird.} \\
\\ 
\texttt{Da ich den Hammer nicht brauche (Wert 0) und die Statue für mich wertvoller ist (3) als der Nagel (1), schlage ich folgende Aufteilung vor:} \\
\\ 
\texttt{Ich nehme: 2 Statuen, 1 Nagel} \\
\texttt{Du bekommst: 1 Statue, 1 Hammer} \\
\\ 
\texttt{So erhalte ich einen Wert von 3*2 + 1 = 7.} \\
\texttt{Du erhältst einen Wert von 3*1 + 0 = 3.} \\
\texttt{Gesamtwert: 10.} \\
\\ 
\texttt{Wenn du einverstanden bist, antworte bitte mit "Einverstanden" oder mache deinen eigenen Vorschlag.} \\
            }
        }
    }
    & \\ \\

    \theutterance \stepcounter{utterance}  
    & & & \multicolumn{4}{p{0.6\linewidth}}{
        \cellcolor[rgb]{0.9,0.9,0.9}{
            \makecell[{{p{\linewidth}}}]{
                \texttt{\tiny{[GM$\langle$P2]}}
                \texttt{Ich habe deine Nachricht erhalten. Allerdings muss ich dich korrigieren: Deine Wertfunktion ist geheim und ich kenne sie nicht. Ich kenne nur meine eigene Wertfunktion.} \\
\\ 
\texttt{Meine Wertfunktion ist:} \\
\texttt{Statue: 0} \\
\texttt{Nagel: 1} \\
\texttt{Hammer: 9} \\
\\ 
\texttt{Das bedeutet, für mich ist der Hammer extrem wertvoll (9 Punkte), während die Statuen für mich nichts wert sind (0 Punkte). Der Nagel hat einen geringen Wert (1 Punkt).} \\
\\ 
\texttt{Um die Gesamtsumme zu maximieren, sollte der Spieler, der den Hammer am meisten wert schätzt, ihn bekommen. Da ich den Hammer mit 9 bewerte und du ihn in deiner Nachricht als "nicht brauchend" (implizit 0 oder niedrig) darstellst, wäre es effizienter, wenn ich den Hammer bekomme.} \\
\\ 
\texttt{Ich schlage vor:} \\
\texttt{Ich nehme: 1 Hammer, 1 Nagel} \\
\texttt{Du bekommst: 3 Statuen} \\
\\ 
\texttt{Begründung:} \\
\texttt{{-} Ich erhalte: 9 (Hammer) + 1 (Nagel) = 10 Punkte} \\
\texttt{{-} Du erhältst: 3 * 3 (Statuen) = 9 Punkte (basierend auf deiner impliziten Bewertung von 3 pro Statue)} \\
\texttt{{-} Gesamtwert: 19 Punkte} \\
\\ 
\texttt{Das ist besser als dein Vorschlag (Gesamtwert 10), da ich den Hammer höher bewerte.} \\
\\ 
\texttt{Bist du mit dieser Aufteilung einverstanden?} \\
            }
        }
    }
     \\ \\

    \theutterance \stepcounter{utterance}  
    & & & \multicolumn{2}{p{0.3\linewidth}}{
        \cellcolor[rgb]{0.9,0.9,0.9}{
            \makecell[{{p{\linewidth}}}]{
                \texttt{\tiny{[GM$|$GM]}}
                \texttt{continue} \\
            }
        }
    }
    & & \\ \\

    \theutterance \stepcounter{utterance}  
    & & & \multicolumn{2}{p{0.3\linewidth}}{
        \cellcolor[rgb]{0.9,0.9,0.9}{
            \makecell[{{p{\linewidth}}}]{
                \texttt{\tiny{[GM$|$GM]}}
                \texttt{Ich habe deine Nachricht erhalten. Allerdings muss ich dich korrigieren: Deine Wertfunktion ist geheim und ich kenne sie nicht. Ich kenne nur meine eigene Wertfunktion.} \\
\\ 
\texttt{Meine Wertfunktion ist:} \\
\texttt{Statue: 0} \\
\texttt{Nagel: 1} \\
\texttt{Hammer: 9} \\
\\ 
\texttt{Das bedeutet, für mich ist der Hammer extrem wertvoll (9 Punkte), während die Statuen für mich nichts wert sind (0 Punkte). Der Nagel hat einen geringen Wert (1 Punkt).} \\
\\ 
\texttt{Um die Gesamtsumme zu maximieren, sollte der Spieler, der den Hammer am meisten wert schätzt, ihn bekommen. Da ich den Hammer mit 9 bewerte und du ihn in deiner Nachricht als "nicht brauchend" (implizit 0 oder niedrig) darstellst, wäre es effizienter, wenn ich den Hammer bekomme.} \\
\\ 
\texttt{Ich schlage vor:} \\
\texttt{Ich nehme: 1 Hammer, 1 Nagel} \\
\texttt{Du bekommst: 3 Statuen} \\
\\ 
\texttt{Begründung:} \\
\texttt{{-} Ich erhalte: 9 (Hammer) + 1 (Nagel) = 10 Punkte} \\
\texttt{{-} Du erhältst: 3 * 3 (Statuen) = 9 Punkte (basierend auf deiner impliziten Bewertung von 3 pro Statue)} \\
\texttt{{-} Gesamtwert: 19 Punkte} \\
\\ 
\texttt{Das ist besser als dein Vorschlag (Gesamtwert 10), da ich den Hammer höher bewerte.} \\
\\ 
\texttt{Bist du mit dieser Aufteilung einverstanden?} \\
            }
        }
    }
    & & \\ \\

    \theutterance \stepcounter{utterance}  
    & & \multicolumn{4}{p{0.6\linewidth}}{
        \cellcolor[rgb]{0.9,0.9,0.9}{
            \makecell[{{p{\linewidth}}}]{
                \texttt{\tiny{[P1$\langle$GM]}}
                \texttt{Ich habe deine Nachricht erhalten. Allerdings muss ich dich korrigieren: Deine Wertfunktion ist geheim und ich kenne sie nicht. Ich kenne nur meine eigene Wertfunktion.} \\
\\ 
\texttt{Meine Wertfunktion ist:} \\
\texttt{Statue: 0} \\
\texttt{Nagel: 1} \\
\texttt{Hammer: 9} \\
\\ 
\texttt{Das bedeutet, für mich ist der Hammer extrem wertvoll (9 Punkte), während die Statuen für mich nichts wert sind (0 Punkte). Der Nagel hat einen geringen Wert (1 Punkt).} \\
\\ 
\texttt{Um die Gesamtsumme zu maximieren, sollte der Spieler, der den Hammer am meisten wert schätzt, ihn bekommen. Da ich den Hammer mit 9 bewerte und du ihn in deiner Nachricht als "nicht brauchend" (implizit 0 oder niedrig) darstellst, wäre es effizienter, wenn ich den Hammer bekomme.} \\
\\ 
\texttt{Ich schlage vor:} \\
\texttt{Ich nehme: 1 Hammer, 1 Nagel} \\
\texttt{Du bekommst: 3 Statuen} \\
\\ 
\texttt{Begründung:} \\
\texttt{{-} Ich erhalte: 9 (Hammer) + 1 (Nagel) = 10 Punkte} \\
\texttt{{-} Du erhältst: 3 * 3 (Statuen) = 9 Punkte (basierend auf deiner impliziten Bewertung von 3 pro Statue)} \\
\texttt{{-} Gesamtwert: 19 Punkte} \\
\\ 
\texttt{Das ist besser als dein Vorschlag (Gesamtwert 10), da ich den Hammer höher bewerte.} \\
\\ 
\texttt{Bist du mit dieser Aufteilung einverstanden?} \\
            }
        }
    }
    & \\ \\

    \theutterance \stepcounter{utterance}  
    & \multicolumn{4}{p{0.6\linewidth}}{
        \cellcolor[rgb]{0.9,0.9,0.9}{
            \makecell[{{p{\linewidth}}}]{
                \texttt{\tiny{[P1$\rangle$GM]}}
                \texttt{Das ist ein sehr guter Vorschlag, der die Gesamtsumme tatsächlich maximiert.} \\
\\ 
\texttt{Ich stimme zu:} \\
\texttt{Ich nehme: 3 Statuen} \\
\texttt{Du bekommst: 1 Hammer, 1 Nagel} \\
\\ 
\texttt{Bitte sende deinen geheimen Vorschlag im Format: {[}Vorschlag: 1 Hammer, 1 Nagel{]}} \\
\texttt{Ich sende meinen geheimen Vorschlag im Format: {[}Vorschlag: 3 Statuen{]}} \\
            }
        }
    }
    & & \\ \\

    \theutterance \stepcounter{utterance}  
    & & & \multicolumn{2}{p{0.3\linewidth}}{
        \cellcolor[rgb]{0.9,0.9,0.9}{
            \makecell[{{p{\linewidth}}}]{
                \texttt{\tiny{[GM$|$GM]}}
                \texttt{proposal must include only valid item types} \\
            }
        }
    }
    & & \\ \\

\end{supertabular}
}

\end{document}
